% Part: sets-functions-relations
% Chapter: infinite
% Section: dedekinds-proof
%
\documentclass[../../../include/open-logic-section]{subfiles}

\begin{document}

\olfileid{sfr}{infinite}{dedekindsproof}

\olsection[Dedekind's ``Proof'']{అనంత సమితి ఉనికి గురించి డెడెకిండ్
``నిరూపణ''}

ఈ అధ్యాయంలో డెడెకిండ్ బీజగణితాల ద్వారా సహజ సంఖ్యలకు సమితి-సిద్ధాంత
వివరణ ఇచ్చాం. \olref[arith][ref]{sec}లో విశ్లేషణ యొక్క
అంకగణితీకరణకు (ఇతర విషయాలతో పాటు) ఉన్న తాత్త్విక ప్రాముఖ్యతను
పరిశీలించాం. ఇప్పుడు ఇక్కడ మనం సాధించిన దాని ప్రాముఖ్యతను పరిశీలించాలి.

\olref[sfr][arith][]{chap} అంతటా సహజ సంఖ్యలు ఇచ్చివున్నాయని తీసుకుని,
వాటిని ఉపయోగించి పూర్ణ, కరణీయ, వాస్తవ సంఖ్యలను స్పష్టంగా నిర్మించాం.
ఈ అధ్యాయంలో సహజ సంఖ్యలకు స్పష్టమైన నిర్మాణాన్ని ఇవ్వలేదు. కేవలం,
\emph{ఏదైనా డెడెకిండ్ అనంత సమితి ఇచ్చినప్పుడు}, $\Nat$ ఎలా
ప్రవర్తించాలని కోరుకుంటామో సరిగ్గా అలాగే ప్రవర్తించే సమితిని
నిర్వచించగలమని చూపాం.

అందువల్ల \emph{సహజ సంఖ్యలు ఏ వస్తువులు} వంటి అధిభౌతిక ప్రశ్నకు
సమాధానం ఇచ్చామని చెప్పలేము. కానీ అది మంచి విషయమే. ఎందుకంటే
\olref[sfr][arith][ref]{sec}లో, $\Real$ను డెడెకిండ్ కోతల సమితిగా
నిర్వచించడాన్ని సౌకర్యవంతమైన నిర్దేశంగా కాకుండా ఒక \emph{కనుగొనికగా}
భావించడం తప్పని నొక్కిచెప్పాం. డెడెకిండ్ కోతలు వాస్తవ సంఖ్యల కీలక
గణిత ధర్మాలకు ఉదాహరణగా నిలుస్తాయనేదే ప్రధాన పరిశీలన. ఇక్కడ కూడా అదే:
\emph{ఏ} డెడెకిండ్ బీజగణితమైనా సహజ సంఖ్యల కీలక గణిత ధర్మాలకు
ఉదాహరణగా నిలుస్తుందనేదే ప్రధాన పరిశీలన. (నిజానికి అన్ని డెడెకిండ్
బీజగణితాలు \emph{సమరూపాలు} అని నిరూపించడం ద్వారా డెడెకిండ్ ఈ విషయాన్ని
మరింత బలంగా స్థాపించాడు (\citeyear[Theorems
132--3]{Dedekind1888}).
అందువల్ల అనేక సమకాలీన ``నిర్మాణవాదులు'' డెడెకిండ్‌ను పూర్వగామిగా
ఉటంకించడంలో ఆశ్చర్యం లేదు.)

ఇంకా, సహజ సంఖ్యల సిద్ధాంతాన్ని సరళమైన సమితి సిద్ధాంతంలో ఎలా
అంతఃస్థాపించాలో చూపాం. ఆ సమితి సిద్ధాంతం ఇప్పటికీ కొంత అనియతమే;
అయినా సహజ సంఖ్యలు ఇచ్చివున్నాయని అది (చూస్తే) ఊహించదు. అందువల్ల
డెడెకిండ్ స్వయంగా ఇలా వివరించిన ప్రతిష్ఠాత్మక ప్రణాళికను సాధించే
మార్గంలో మనం ఉండవచ్చు:
\begin{quote}
	విజ్ఞానంలో నిరూపించగలిగే దేనినీ నిరూపణ లేకుండా నమ్మకూడదు. ఈ
	అవసరం సమంజసంగా కనిపించినా, అత్యంత సరళమైన విజ్ఞానానికి పునాదులు వేసే
	అత్యంత ఇటీవలి పద్ధతుల్లో కూడా---అంటే సంఖ్యల సిద్ధాంతంతో వ్యవహరించే
	తర్కశాస్త్ర భాగంలో కూడా---ఇది నెరవేరిందని నేను భావించలేను.
	అంకగణితాన్ని (బీజగణితం, విశ్లేషణను) కేవలం తర్కశాస్త్రంలోని భాగంగా
	చెప్పడంలో, సంఖ్య అనే భావన స్థలం, కాలం గురించిన భావనలు లేదా
	అంతర్బుద్ధుల నుంచి పూర్తిగా స్వతంత్రమని; దానిని శుద్ధ ఆలోచనా
	నియమాల తక్షణ ఫలితంగా భావిస్తానని సూచిస్తున్నాను.
	\citep[preface]{Dedekind1888}
\end{quote}
డెడెకిండ్ సాహసోపేత ఆలోచన ఇది. కేవలం (సరళ) సమితి సిద్ధాంతాన్ని
ఉపయోగించి సహజ సంఖ్యలను ఎలా నిర్మించాలో ఇప్పుడే చూపాం.
\olref[sfr][arith][]{chap}లో, సహజ సంఖ్యలు, కొంత సమితి సిద్ధాంతం
ఇచ్చినప్పుడు వాస్తవ సంఖ్యలను ఎలా నిర్మించాలో చూశాం. అందువల్ల బహుశా
``అంకగణితం (బీజగణితం, విశ్లేషణ)'' అనేది ``కేవలం తర్కశాస్త్రంలోని
భాగమే'' అవుతుంది (``తర్కశాస్త్రం'' అనే పదానికి డెడెకిండ్ ఇచ్చిన
విస్తృత అర్థంలో).

ఆలోచన అదే. కానీ ఒక్క క్షణం ఆగండి. సహజ సంఖ్యలకు మన ప్రతిరూపమైన
డెడెకిండ్ బీజగణితం నిర్మాణం, ఒక డెడెకిండ్ అనంత సమితి ఉనికిపై ఆధారపడి
ఉంది. (\olref[sfr][infinite][dedekind]{thm:DedekindInfiniteAlgebra}ను
మళ్లీ చూడండి.) డెడెకిండ్ అనంత సమితి ఉనికిని ``తర్కశాస్త్రం'' లేదా
``శుద్ధ ఆలోచనా నియమాల'' ద్వారా స్థాపించలేకపోతే, ప్రణాళిక అక్కడే
ఆగిపోతుంది.

అయితే డెడెకిండ్ అనంత సమితి ఉనికిని ``శుద్ధ ఆలోచనా నియమాల'' ద్వారా
స్థాపించగలమా? డెడెకిండ్ ప్రయత్నం ఇదిగో:
\begin{quote}
  నా ఆలోచనల స్వంత రాజ్యం, అంటే నా ఆలోచనకు వస్తువులు కాగల అన్ని
  విషయాల సమగ్రత $S$, అనంతమైనది. ఎందుకంటే $s$, $S$లోని ఒక మూలకాన్ని
  సూచిస్తే, $s$ నా ఆలోచనకు వస్తువు కాగలదనే ఆలోచన $s'$ కూడా
  $S$లోని ఒక మూలకం. దీనిని $s$ అనే మూలకపు ప్రతిబింబం $\phi(s)$గా
  పరిగణిస్తే, \dots~$S$ [డెడెకిండ్] అనంతమైనది; నిరూపించాల్సింది ఇదే.
	\citep[\S66]{Dedekind1888}
\end{quote}
అత్యంత కచ్చితమైన గణిత నిరూపణలతో ఎక్కువగా నిండిన ఒక పుస్తకం మధ్యలో
ఇలాంటి విషయం కనిపించడం నిజంగా ఆశ్చర్యకరం. రెండు వ్యాఖ్యలు చేయాలి.

మొదట: ఈ ``నిరూపణ''కు ఇప్పుడు మనం ``గణితపరమైనది''గా గుర్తించే లక్షణం
దాదాపు లేదు. ఇది మానసిక వస్తువుల (ఆలోచనల) గురించి, అవి కూడా కేవలం
\emph{సంభవనీయమైన} వాటి గురించి మాట్లాడుతుంది.

రెండవది: కనీసం మనం డెడెకిండ్ బీజగణితాలను వివరించిన విధానంలో, ఈ
``నిరూపణ''కు ఒక సూటి సాంకేతిక లోపం ఉంది. డెడెకిండ్ వాదన విజయవంతమైతే,
అనంత విషయాలు (ప్రత్యేకంగా అనంత ఆలోచనలు) ఉన్నాయని మాత్రమే స్థాపిస్తుంది.
కానీ $S$ను బహువచనంలో కేవలం \emph{కొన్ని విషయాలు}గా భావించడానికి బదులుగా,
అనంత !!{element}s గల ఒకే \emph{సమితి}గా $S$ను పరిగణించడానికి కూడా డెడెకిండ్
ఒక కారణాన్ని ఇవ్వాలి.

ఇక్కడ ఖాళీ ఉందని డెడెకిండ్ చూడకపోవడం, ఆయన ``సమగ్రత'' అనే పదాన్ని
వాడిన తీరు \emph{మనం} ``సమితి'' అనే పదాన్ని వాడే తీరుతో కచ్చితంగా
సరిపోలదని సూచించవచ్చు.\footnote{నిజానికి అది సరిపోలలేదని భావించడానికి
మనకు ఇతర కారణాలూ ఉన్నాయి; \citet[p.~23]{Potter2004} చూడండి.}
కానీ ఇది పెద్ద ఆశ్చర్యం కాదు. గత రెండు అధ్యాయాల్లో మనం అనుసరించిన
ప్రణాళిక---సహజ సంఖ్యల ``నిర్మాణం'', వాటి నుంచి పూర్ణ, వాస్తవ,
కరణీయ సంఖ్యల ``నిర్మాణం''---అంతా సరళ పద్ధతిలోనే జరిగింది. కచ్చితంగా
ఏ సమితులు ఎప్పుడు ఉనికిలో ఉంటాయనే అనేక సాధారణ సూత్రాలను పెట్టకుండా,
అవసరమైనప్పుడల్లా ఈ సమితిని లేదా ఆ సమితిని ఉపయోగించుకున్నాం. కానీ
\emph{కొన్ని} సాధారణ సూత్రాలు అవసరమని మనకు తెలుసు; లేకపోతే రసెల్
వైరుధ్యంలో పడతాం.

ఇక మన సరళత్వాన్ని అధిగమించాల్సిన సమయం వచ్చింది.

\end{document}
